% Generated from paper/full_draft. Edit the Markdown source or migration script.

\section{Introduction}
\label{sec:02-introduction:introduction}

Knowledge-augmented agents improve language-model responses by conditioning generation on external evidence. In many systems this evidence is supplied by a retrieval-augmented generation (RAG) layer, but the underlying control problem is broader: an agent must decide which pieces of structured domain knowledge should enter its limited context. If relevant evidence is not selected, the generator has little chance to recover; if too much evidence is selected, latency, context cost, and distracting noise increase. Practical systems must repeatedly decide how much evidence to retrieve, which route to trust, and when a smaller context is safe.

Dense retrieval is a strong baseline for this problem because it can recover semantically related passages even when the query and evidence do not share the same surface form. However, dense-only retrieval remains a fixed route. It does not explicitly decide when lexical matching is safer, when local cluster structure is informative, when global dense retrieval should remain as a fallback, or when the final retrieved context can be compacted. These questions become more important in vertical-domain knowledge settings, where corpora are often large, terminology-heavy, and shaped by repeated workflows, entities, and user intents.

This paper studies the following hypothesis: vertical-domain knowledge data often exhibits a piecewise relevance structure. Relevant evidence is not uniformly distributed across the embedding space. Instead, it tends to form local regions induced by domain concepts, lexical anchors, task workflows, and user behavior. Pure dense retrieval performs local semantic neighborhood search in one representation space, but it may not fully capture lexical constraints, cluster-level routing, or user-specific relevance. At the same time, geometry alone is not enough: a cluster route that prunes too early can miss the correct evidence, and dense retrieval must remain available as a recall floor.

We propose IntentWeight, a feedback-guided adaptive evidence selection controller motivated by a piecewise relevance-manifold assumption for vertical-domain data. IntentWeight does not replace dense retrieval with a single alternative retriever. Instead, in our retrieval-augmented QA implementation, it builds a multi-route retrieval surface including dense retrieval, BM25 lexical recall, and cluster-local retrieval. Fixed KMeans/MiniBatchKMeans clusters provide stable arms for a LinUCB policy. The policy observes query and route features, selects cluster-local routes, and is updated by trust-weighted simulated feedback. A confidence-based final context policy then decides whether to compact the retrieved context or keep a denser fallback.

Although the control problem appears in memory, graph, tree, and retrieval-backed agents, our empirical validation instantiates it in retrieval-augmented question answering. We therefore use the broader agent framing as motivation and keep the demonstrated claims tied to retrieval-backed evidence selection.

This design also separates retrieval ranking from final context control. A late reranker can improve the ordering of a candidate pool, and a sentence compressor can remove redundant text from selected evidence. These components are useful, but they do not by themselves decide which retrieval route should be trusted, when dense fallback is necessary, or how much final evidence risk a query can tolerate. IntentWeight is positioned as the route-and-budget controller in this stack, not as a replacement for dense retrieval, reranking, or context compression.

The same distinction separates three cost layers that are often conflated in RAG experiments: the number of source candidates considered during retrieval, the rate at which global dense retrieval is invoked, and the final number of retrieved context tokens sent to the generator. Our main efficiency claim uses the third layer. Earlier candidate-count reductions are useful retrieval-stage diagnostics, but they are not evidence of lower LLM context cost unless the final context itself is reduced. Because these retrieved chunks enter the LLM generator as input tokens, each percentage point of evidence-context reduction translates directly into a proportional per-query inference-cost reduction, a recurring saving that scales with deployment query volume.

We evaluate IntentWeight on multiple datasets and use LoTTE technology/search as the main large-scale vertical-domain evidence benchmark. On LoTTE, we scale from 100k to 638k corpus chunks and compare against dense-only retrieval using \nolinkurl{sentence-transformers/all-MiniLM-L6-v2} with exact cosine search. Our main cost-quality evidence uses a calibration/test context-budget protocol: the final-context policy is selected on calibration queries and frozen before test evaluation. Under this protocol, calibration-eligible operating points at 100k, 200k, and 638k save 6-18\% final evidence-context tokens, while a 400k diagnostic point shows positive frozen-test behavior but does not satisfy the calibration eligibility gate. Across these scales, IntentWeight avoids the larger $\mathrm{Hit@10}$ losses observed under dense-only adaptive truncation, although strict seed-level non-inferiority remains scale-dependent. A conservative confidence-only policy provides a stable baseline, reducing final retrieved context tokens by approximately 4.7-5.3\% across all scales while preserving dense-level query hit. We treat these as bounded operating points rather than universal or statistically significant dominance claims.

We further evaluate three reviewer-facing concerns. First, sentence-level MMR over dense top-10 preserves dense chunk-support while reducing selected sentence tokens, showing that simple final-context compression is a strong baseline and should be applied uniformly when used. Second, a cross-encoder reranker over dense top-50 improves full top-10 support metrics but selects longer chunks on average; under the same per-query token budgets used by IntentWeight, it does not uniformly dominate the calibrated controller. Third, a LoTTE science/search replication shows that the fixed top-10 IntentWeight ranking improves query-level $\mathrm{Hit@10}$ at both 20k and 100k corpus scales, while also showing that aggressive context compression must be domain-calibrated. Finally, a hard-case recovery experiment shows that simulated arm-level feedback can repair a meaningful fraction of budget-induced tail failures in post-feedback retry. This recovery result supports the adaptive feedback mechanism, but it is not a claim of first-pass IID improvement on all future queries.

The contributions of this paper are:

\begin{enumerate}
\item We formulate evidence selection over structured vertical-domain data as an adaptive route-control problem motivated by local relevance structure rather than a fixed retriever selection problem.
\item We state a bounded piecewise relevance-manifold hypothesis and test its diagnostic implications with local geometry measurements across LoTTE domains and scales.
\item We introduce IntentWeight, a feedback-adaptive multi-route controller combining dense retrieval, BM25 lexical recall, cluster-local retrieval, trust-weighted LinUCB route learning, and confidence-based final context compaction in a retrieval-augmented QA implementation.
\item We provide large-scale LoTTE evidence that calibrated context-budget policies reduce the LLM input tokens consumed by retrieved evidence context while avoiding the larger $\mathrm{Hit@10}$ losses of dense-only adaptive truncation under bounded operating points.
\item We add strong post-retrieval baselines: dense sentence-MMR compression, a compressor-normalized dense-versus-IntentWeight comparison, and a cross-encoder reranker same-budget check. These baselines position IntentWeight as a route-and-budget controller that is complementary to compression and reranking.
\item We report paired query-level statistics and calibration eligibility to separate retrieval-quality non-inferiority from token-cost reduction.
\item We replicate the ranking-side result on a second LoTTE domain and document that context-compression strength requires domain calibration.
\item We show that simulated feedback can act as a controlled recovery mechanism for tail queries harmed by aggressive context compression.
\item We document limitation cases where dataset structure, weak labels, duplicate evidence, sparse ground truth, or complete-evidence requirements reduce the benefit of adaptive routing.
\end{enumerate}

The resulting claim is intentionally bounded. IntentWeight is not presented as a universal replacement for dense retrieval. It is a feedback-driven controller that uses dense retrieval as a recall floor and learns when route confidence is strong enough to reduce the final context budget or trigger a safer post-feedback recovery path. Reranking and sentence compression remain compatible downstream layers rather than competing explanations to hide.
